% Generated draft: requires mathematical and source review.
\documentclass[11pt]{article}
\usepackage[margin=25mm]{geometry}
\usepackage{fontspec}
\setmainfont{Latin Modern Roman}
\usepackage{amsmath,amssymb,mathtools,unicode-math}
\setmathfont{Latin Modern Math}
\usepackage{enumitem,needspace,xurl,hyperref}
\setlength{\parindent}{0pt}
\setlength{\parskip}{5pt}
\setlength{\emergencystretch}{3em}
\newcommand{\N}{\mathbb{N}}
\newcommand{\Z}{\mathbb{Z}}
\newcommand{\Q}{\mathbb{Q}}
\newcommand{\R}{\mathbb{R}}
\newcommand{\C}{\mathbb{C}}
\newcommand{\F}{\mathbb{F}}
\newcommand{\A}{\mathbb{A}}
\newcommand{\PP}{\mathbb{P}}
\newcommand{\E}{\mathbb{E}}
\newcommand{\eps}{\varepsilon}
\newcommand{\dd}{\,\mathrm{d}}
\begin{document}
\section*{UnsolvedMath ID: 30003905}
\textbf{Source problem code:} catalogue entry
\subsection{Definitions and mathematical statement}
For a number field \(K\), let \(\operatorname{Cl}(K)\) denote its ideal class group and let \(\mathcal O_K^{\times}\) denote its unit group. For a prime \(\ell\) and a finite abelian group \(A\), written additively, its \(\ell\)-primary part is
\[
A[\ell^{\infty}]=\{a\in A:\ell^n a=0\text{ for some integer }n\geq 0\}.
\]
The supplied record does not specify which primes \(\ell\) are admissible or what definition of relative class group is intended.

The source proposes the following conjectural comparison: after incorporating the contribution of unit groups, the joint distribution over number fields of the relative class-group \(\ell\)-part and two additional class-group invariants agrees with the distribution predicted by a random symplectic-similitude model. Its intended mathematical form can therefore be recorded only schematically as
\[
\mu_{\mathrm{arithmetic}}=\mu_{\mathrm{symplectic\text{-}similitude}},
\]
where both sides are intended to describe the same triple of invariants, with the unit-group contribution incorporated as required by the source.

This is not yet a fully specified equality of probability measures: the record supplies neither the family of number fields and underlying extensions, nor their ordering or probability law, nor the relative class-group convention, nor the two accompanying invariants. It also does not define the unit-group contribution, the parameters and probability law of the random symplectic-similitude model, or the procedure extracting the triple from that model. These omissions prevent a precise quantified conjecture from being recovered from the supplied record alone.
\subsection{Short English statement}
The proposed conjecture compares the joint distribution of three class-group invariants over number fields with a random symplectic-similitude prediction, taking unit groups into account. The supplied record leaves the field family, two of the invariants, and the model unspecified, so it determines only a schematic statement requiring mathematical review.
\subsection{Sources}
\begin{itemize}
\item Local UnsolvedMath registry, supplied source record with permanent ID \(30003905\), title \emph{Symplectic Models for Relative Class Group Distributions}, source problem code \texttt{catalogue entry}, in the Oberwolfach Reports Open Problems collection.
\item Supplied report citation: \emph{Explicit Methods in Number Theory (2019)}, Oberwolfach Reports, pp.\ 2033--2076, DOI \texttt{10.4171/owr/2018/34}. Supplied source link: \url{https://doi.org/10.4171/owr/2018/34}. The registry gives proposed year 2018 and names no proposer.
\item Supplied literature source, labeled ``OWR report'': \url{https://ems.press/content/serial-article-files/46754}.
\item Literature assessment dated exactly \texttt{2026-08-21}: the supplied curation labels the problem open and states that the joint distribution over number fields is explicitly conjectural, with function-field verification and numerical evidence in the report but no proof for number fields. This is supplied dated curation, not an independently verified literature assessment; its time-sensitive status and cited claims require independent verification.
\item The registry labels the statement \texttt{corrected\_verified} and says that notation, typography, and sentence structure were repaired while preserving mathematical content. This supplied editorial qualification does not resolve the missing mathematical specifications or establish mathematical correctness.
\end{itemize}
\end{document}
